\section{Verification by sampling}
\label{sec:13.5}
Verification from outside is partly cryptographic, and it works by sampling. The operator commits in advance to the deployed weights, and for each period to the full log of production inputs and outputs. After each commitment, entries are drawn by randomness the operator can't predict, such as a public beacon, and the operator proves those entries came from the committed weights without revealing the weights or the inputs. If a fraction f of outputs came from another model, k sampled proofs catch it with probability 1 − (1 − f)\textsuperscript{k}; substitution on one output in a hundred is caught about 95 percent of the time at 300 proofs. That catches silent restoration, retraining and substitution, which change every output. Sampling misses substitution on a single consequential output, so an adverse party holds a second right: to demand proofs for specific logged outputs once their consequences are known. In a targeting pipeline those outputs are few and identifiable afterwards, because an airstrike happened. The right requires the operator to retain committed logs for a stated period and to link consequential actions to log entries. Each entry records the class the case was sorted into, and the targeted right covers auto-proceeded entries in particular, since that is where a consequential output is likeliest to have had no person look at it. Outside probes (whether a refusal tracks its reason, how far the system discounts a claim by its source's stake) are proofs on inputs an outsider chose, a handful at a time. The committed log's length is the generation rate, and the class recorded in each entry makes the partition's proportions provable the same way. Neither reveals the targets, and neither attests that the log is complete.

Sampling proves that logged outputs came from the committed model and says nothing about outputs never logged; an operator can run an uncommitted model beside the committed one, which is routing. And a proof shows f(x) = y for the x committed, and is silent on whether x was all the evidence. Proving even one frontier-scale forward pass is not yet practical, and production inference adds nondeterministic numerics and sampling randomness that a proof would have to fix in advance. What the proofs would attest is stated now, so that a deployment can be held to it once they can be produced. The residual (whether the operator routed around the bearer, and what lay behind compartments it wasn't cleared into) is the residual with a human analyst. That isn't a reassurance on its own, since the officer was a human analyst. What distinguished his failure was tempo, and the review ratio is what addresses it.
